\section{Assumptions}
\label{sec:assumptions}

The analysis rests on \emph{two} assumptions: a single global behavioural
assumption \Cref{ass:descent} on the per-step \emph{correct-direction}
alignment, and an architectural assumption \Cref{ass:bounded_architecture}
consolidating deterministic bounds on the unembedding, the value vectors,
and the attention projections (including an explicit bounded-score regime).
The behavioural assumption carries \emph{two} interpretable knobs: a
\emph{rate} $p$ (how often a step points towards the answer) and a
\emph{magnitude} $c$ (how much). These multiply to the net per-step drift
$\snet\driftc\rho_0$ of \Cref{lem:signal_floor}; convexity of the loss
motivates that both $\driftc>0$ and $p>\tfrac12$ are plausible
(\Cref{rem:descent_naming}) but does \emph{not} derive their values.

\subsection{The global alignment assumption}
\label{subsec:descent}

Recall from \Cref{def:descent_indicator} the unit correct-token row
$\rowhat=W_U^{\corr}/\norm{W_U^{\corr}}_2$, the per-step
\emph{correct-direction projection}
$\proj_k\coloneqq\inner{V_k}{\rowhat}$ of the value vector $V_k$, and the
sign indicator $\descent_k\coloneqq\operatorname{sign}(\proj_k)\in\{-1,+1\}$.

\begin{assumption}[Global correct-direction alignment: rate $p$ and magnitude $\driftc$]
\label{ass:descent}
There are constants $p\in[0,1]$ and $\driftc\in(0,1]$ such that, for every
step $k\ge1$, the correct-direction projection $\proj_k$ of
\Cref{def:descent_indicator} satisfies all three of:
\begin{enumerate}[label=(\alph*),leftmargin=2em,itemsep=0.2em]
\item \emph{Sign rate.} The sign indicator obeys
       \begin{align}\label{eq:descent_condition}
          \Pr\bigl[\, \descent_k = +1 \,\bigm|\, \Fcal_{k-1} \,\bigr] \;=\; p,
       \end{align}
       on \emph{every} step, with no dependence on the current loss value,
       on the region of state space, or on the time index;
\item \emph{Magnitude floor.} the conditional mean magnitude is floored,
       \begin{align}\label{eq:magnitude_floor}
          \E\bigl[\, \abs{\proj_k} \,\bigm|\, \Fcal_{k-1} \,\bigr]
          \;\ge\; \driftc\,\rho_0,
       \end{align}
\item \emph{Sign/magnitude conditional independence.} given $\Fcal_{k-1}$,
       the sign $\descent_k$ and the magnitude $\abs{\proj_k}$ are
       conditionally independent;
\item \emph{Drift direction (systematic update points at the answer).} the
       conditional mean of the value vector is aligned with the correct
       row, i.e.\ there is a $\Fcal_{k-1}$-measurable scalar $\beta_k$ with
       \begin{align}\label{eq:drift_direction}
          \E[V_k\mid\Fcal_{k-1}] \;=\; \beta_k\,\rowhat,
          \qquad
          \beta_k \;=\; \E[\proj_k\mid\Fcal_{k-1}],
       \end{align}
       so that the only systematic direction of a step is towards
       $W_U^{\corr}$ and everything else is mean-zero noise
       (\Cref{ass:bounded_architecture}(5)).
\end{enumerate}
Define the \emph{net alignment rate}
\begin{align}\label{eq:delta_def}
   \snet \;\coloneqq\; 2p - 1 \;\in\; [-1, 1].
\end{align}
\end{assumption}

\begin{remark}[Globality, the two knobs, and the rigorous net drift]
\label{rem:descent_globality}
\Cref{ass:descent} is the single behavioural hypothesis of the paper.
\emph{(i) Global.} \Crefrange{eq:descent_condition}{eq:magnitude_floor}
hold for all $k$ with the \emph{same} constants $p,\driftc$, with no
threshold $\Lstar$ gating the rate and no privileged sublevel set of
$\loss$: there is no loss-basin, no ``snowball'' region, and no
stratification of the indicator by the loss value. \emph{(ii) No initial
condition.} $x_0$ is arbitrary on $\Sphere$ (the antipodal worst case
included), because $x_0$ washes out of the un-retracted iterate at rate
$w_{T,0}=O(1/T)$ (\Cref{eq:running_average_teaser}) and the
isotropic-noise non-degeneracy of \Cref{ass:bounded_architecture}(5)
prevents the trajectory from being trapped at the measure-zero set
collinear with $W_U^{\corr}$ (where $\proj_k$ would degenerate). \emph{(iii)
Two knobs, stated on the correct direction.} Clause~(a) is a \emph{sign}
condition on the realised projection onto $\rowhat$; clause~(b) is its
\emph{magnitude}. These are the two interpretable knobs --- rate and
magnitude --- and they multiply to the net drift:

\smallskip\noindent\textbf{The net per-step drift, derived.} By
clause~(c), conditional on $\Fcal_{k-1}$ the sign $\descent_k$ and the
magnitude $\abs{\proj_k}$ are independent, so since
$\proj_k=\descent_k\abs{\proj_k}$ (with $\descent_k\in\{-1,+1\}$),
\begin{align}\label{eq:net_drift}
   \E[\proj_k\mid\Fcal_{k-1}]
   \;=\; \E[\descent_k\mid\Fcal_{k-1}]\cdot\E[\abs{\proj_k}\mid\Fcal_{k-1}]
   \;=\; \snet\,\E[\abs{\proj_k}\mid\Fcal_{k-1}],
\end{align}
using $\E[\descent_k\mid\Fcal_{k-1}]=p-(1-p)=\snet$ from
\Cref{eq:descent_condition}. \Cref{eq:net_drift} is an \emph{identity}
(valid for either sign of $\snet$); combined with the bounds
$\driftc\rho_0\le\E[\abs{\proj_k}\mid\Fcal_{k-1}]\le M$ (the floor
\Cref{eq:magnitude_floor} and $\abs{\proj_k}\le\norm{V_k}_2\le M$ from
\Cref{ass:bounded_architecture}(2)) it gives, whenever $\snet\ge0$, the
correct-direction drift floor
\begin{align}\label{eq:net_drift_floor}
   \E[\proj_k\mid\Fcal_{k-1}] \;\ge\; \snet\,\driftc\,\rho_0,
\end{align}
and the two-sided control $\abs{\E[\proj_k\mid\Fcal_{k-1}]}\le\abs{\snet}M$
in general. This is the rigorous net drift consumed by
\Cref{lem:signal_floor}; crucially it is obtained \emph{directly} from the
correct-direction rate and magnitude, with \emph{no} transitivity through
the loss-descent direction $-\rgrad\loss$. The super-critical hypothesis
$\snet\ge\sqrt2\,\critrate$ of \Cref{thm:R1} is the substantive question;
$\snet>0$ is its \emph{consequence}, never a separate assumption.
\end{remark}

\begin{remark}[On clause (c) and the anti-aligned tail]
\label{rem:descent_independence}
Clause~(c) is the device that controls the \emph{anti-aligned} steps
($\descent_k=-1$): without it, a step that points away from $\rowhat$ could
carry an arbitrarily large magnitude and overwhelm the favourable steps, so
the net drift \Cref{eq:net_drift} would not factor. Conditional
independence of sign and magnitude is the clean sufficient condition and
makes \Cref{eq:net_drift} an exact identity. It is satisfied, for instance,
whenever the conditional law of $\proj_k$ given $\Fcal_{k-1}$ is a
sign-mixture of a fixed nonnegative magnitude law (a Rademacher-type model
with rate $p$), or more generally whenever the conditional density of
$\proj_k$ is the product of a sign mass and an even magnitude profile.
A strictly symmetric two-sided floor $\abs{\proj_k}\ge\driftc\rho_0$ a.s.\
on both sign branches is an alternative that yields the same conclusion; we
state the conditional-mean form \Cref{eq:magnitude_floor} as it is the
weakest hypothesis that closes \Cref{lem:signal_floor}. No anti-aligned
step is left with unbounded magnitude.

Clause~(d) is the companion structural hypothesis: it says the
\emph{systematic} (conditional-mean) part of each reasoning step points at
the answer $W_U^{\corr}$, while all other variability is mean-zero noise.
It is consistent with clause~(b) --- projecting \Cref{eq:drift_direction}
onto $\rowhat$ gives $\E[\proj_k\mid\Fcal_{k-1}]=\beta_k\norm{\rowhat}_2^2=\beta_k$
since $\norm{\rowhat}_2=1$ --- and it is what lets the incoherence transfer
the drift to a small bound on the \emph{incorrect} rows in
\Cref{lem:incorrect_max}: for $a\neq\corr$,
$\abs{\inner{W_U^a}{\E[V_k\mid\Fcal_{k-1}]}}=\abs{\beta_k}\,\abs{\inner{W_U^a}{\rowhat}}\le M\,\incoh_0 R_U$,
using $\abs{\beta_k}\le M$ and
$\abs{\inner{W_U^a}{\rowhat}}=\abs{\inner{W_U^a}{W_U^{\corr}}}/\norm{W_U^{\corr}}_2\le\incoh_0 R_U$.
Like the rest of \Cref{ass:descent} it is \emph{global} (every $k$) and
imposes no region, basin, or initial condition.
\end{remark}

\subsection{Bounded architectural quantities}
\label{subsec:bounded_architecture}

\begin{assumption}[Bounded architecture and bounded-score regime]
\label{ass:bounded_architecture}
The model architecture and its trained parameters satisfy:
\begin{enumerate}[label=(\arabic*),leftmargin=2em,itemsep=0.2em]
\item \emph{Unembedding $W_U$.} The row-incoherence of
       \Cref{def:incoherence} satisfies $\incoh(W_U)\le\incoh_0$, and the
       row norms satisfy $0<\rho_0\le\norm{W_U^a}_2\le R_U<\infty$ for
       every token $a\in\Vocab$.
\item \emph{Value vectors.} $\norm{V_t}_2\le M$ almost surely for every
       $t\ge1$, with constant $M>0$.
\item \emph{Query/key projections.}
       $\norm{W_Q}_{\mathrm{op}}\le W_{\!QK}$ and
       $\norm{W_K}_{\mathrm{op}}\le W_{\!QK}$ deterministically, with
       constant $W_{\!QK}>0$.
\item \emph{Bounded-score regime.} The scaled dot-product attention
       scores satisfy $\abs{\inner{q}{k_t}}\le S$ for every $t$, with a
       constant $S=O(1)$ in $d$.
\item \emph{Isotropic noise (two-sided, scale $1/d$).} Write
       $V_t=\E[V_t\mid\Fcal_{t-1}]+\xi_t$ with $\xi_t$ the mean-zero noise
       component. Conditional on $\Fcal_{t-1}$, $\xi_t$ has covariance
       comparable to isotropic at scale $1/d$,
       \begin{align}\label{eq:noise_isotropy}
          \frac{\sigma_{\min}^2}{d}\,I
          \;\preceq\;
          \operatorname{Cov}(\xi_t\mid\Fcal_{t-1})
          \;\preceq\;
          \frac{\sigma_{\max}^2}{d}\,I
          \qquad\text{on } T_{x_{t-1}}\Sphere,
       \end{align}
       for constants $0<\sigma_{\min}\le\sigma_{\max}<\infty$. The lower
       bound is the non-degeneracy used for anti-concentration
       (\Cref{lem:incorrect_max_lower}) and the anti-pole washout; the
       upper bound is the per-direction fluctuation scale fed to
       \Cref{lem:azuma}. Both express that the noise is high-dimensional:
       a fixed-direction projection carries an $O(1/d)$ fraction of its
       energy (\Cref{lem:orthogonality}). For the anti-concentration lower
       bound of \Cref{lem:incorrect_max_lower} only, we use the
       \emph{exactly isotropic Gaussian} specialisation of this clause:
       conditional on $\Fcal_{t-1}$,
       \begin{align}\label{eq:noise_gaussian}
          \xi_t \;\bigm|\; \Fcal_{t-1}
          \;\sim\; \mathcal N\!\bigl(0,\ \tau_t^2\,I_d\bigr),
          \qquad
          \sigma_{\min}^2/d \;\le\; \tau_t^2 \;\le\; \sigma_{\max}^2/d,
       \end{align}
       with $\tau_t^2$ an $\Fcal_{t-1}$-measurable scalar.
       \Cref{eq:noise_gaussian} sharpens \Cref{eq:noise_isotropy} (a scalar
       covariance is a special case of one comparable to isotropic) by
       fixing the per-step covariance to be a \emph{scalar} multiple of the
       identity and the law to be Gaussian; the tangential restriction of
       \Cref{eq:noise_isotropy} is dropped here because the radial component
       of a step leaves the incorrect-logit fluctuations unchanged at
       leading order and is removed by the retraction
       (\Cref{lem:retraction_stability}). It is used \emph{only} by
       \Cref{lem:incorrect_max_lower} (see \Cref{rem:gaussian_noise_scope}).
\end{enumerate}
\end{assumption}

\begin{remark}[Realism, the (non-load-bearing) incoherence gap, and when $S=O(1)$ fails]
\label{rem:architecture_realism}
Clause~(1) formalises ``token embeddings are not too colinear''; for
$W_U$ with rows near-uniform on the sphere the typical incoherence is
$\incoh(W_U)=O(\sqrt{\log|\Vocab|/d})\ll1$ by standard concentration
(\cite{vershynin2018}, Chapter~3), and the canonical unit-norm convention
$\rho_0=R_U=1$ makes $\incoh_0\ll1$. The incoherence $\incoh_0$ is consumed
in exactly two places: the deterministic drift term $\incoh_0 R_U M$ of
\Cref{lem:incorrect_max} (orthogonality of the expected iterate to the
incorrect rows) and the pairwise decorrelation $\le\incoh_0\sigma_T^2$ of
\Cref{lem:incorrect_max_lower}. We do \emph{not} require any
\emph{incoherence gap} $\incoh_0<\rho_0^2/R_U^2$: the per-step magnitude
$\driftc$ is now \emph{assumed} positive (\Cref{ass:descent}), not derived
from a gap, so no gap condition is load-bearing. (The gap is the condition
under which the \emph{motivating} gradient computation of
\Cref{rem:descent_naming} would be sign-definite; it appears there as
context only.) Clauses~(2)--(3) are standard for attention-trajectory
analyses and hold for any pretrained transformer with LayerNorm-bounded
activations and finite weights (\cite{vaswani2017attention}). Clause~(4) is
the explicit \emph{bounded-score regime}: under (2)--(3) and
$\norm{x_t}_2=\sqrt d$, Cauchy--Schwarz gives
$\abs{\inner{q}{k_t}}\le W_{\!QK}^2\,d/\sqrt{d_k}$ in raw units, so
$S=O(1)$ is a genuine hypothesis on the normalisation convention rather
than an automatic consequence. \emph{When it fails:} under raw
bounded-entry parametrisations with $\norm{x}_2=\Theta(\sqrt d)$ in
absolute units, $S$ can grow like $\sqrt d$, in which case the
quadratic-variation bound $S_T\le e^{2S}/T$ of \Cref{lem:quad_variation}
becomes vacuous and the headline horizon scaling is lost; we therefore
state $S=O(1)$ as a named hypothesis and flag this regime explicitly.
Clause~(5) is the non-degeneracy that, together with the $x_0$-washout of
\Cref{rem:descent_globality}, keeps the trajectory off the measure-zero
antipodal/collinear set and is absorbed into the high-probability budgets
of \Cref{thm:R1,thm:R2,thm:R3prime} (no initial-angle condition is
imposed; see \Cref{rem:descent_globality}).
\end{remark}

\begin{remark}[Scope of the isotropic-Gaussian specialisation \Cref{eq:noise_gaussian}]
\label{rem:gaussian_noise_scope}
The Gaussian form \Cref{eq:noise_gaussian} is invoked in \emph{exactly one}
place: the anti-concentration lower bound \Cref{lem:incorrect_max_lower}
(the failure branch of \Cref{thm:R1}), where it makes the centred incorrect
projections $\inner{W_U^a}{\tilde x_T-\E\tilde x_T}$ \emph{jointly Gaussian}
across $a$, so that Sudakov minoration and the Borell--TIS Gaussian
concentration inequality (\cite{vershynin2018}, Theorems~7.4.1 and~5.2.3)
apply directly and the two-sided transition is unconditional in $\incoh_0$.
\emph{No other result uses Gaussianity.} The success branch of
\Cref{thm:R1} (\Cref{lem:incorrect_max,lem:signal_floor} via the
Azuma/Freedman bound \Cref{lem:azuma}) and the critical-window theorem
\Cref{thm:R2} (via the Hall--Heyde martingale Berry--Esseen,
\Cref{thm:hallheyde_be}) use only the second-moment information in
\Cref{eq:noise_isotropy} --- the conditional-variance bound
$\inner{W_U^a}{\operatorname{Cov}(\xi_t\mid\Fcal_{t-1})W_U^a}\le\sigma_{\max}^2\norm{W_U^a}_2^2/d$
and the a.s.\ increment bound $\norm{\xi_t}_2\le2M$ --- and never the
Gaussian law: indeed \Cref{thm:R2} \emph{derives} asymptotic normality of
the margin from the sub-Gaussian increments via the martingale CLT, which
would be circular if Gaussianity were assumed there. The scalar
(as opposed to merely comparable-to-isotropic) covariance in
\Cref{eq:noise_gaussian} is what reduces the pairwise correlation of the
$Z_a$ to the normalised Gram entry $\incoh(W_U)$ of \Cref{def:incoherence};
it is the natural idealisation of the high-dimensional noise already
expressed by \Cref{eq:noise_isotropy} and is mild for the same reason that
clause is (the noise is direction-agnostic at scale $1/d$).
\end{remark}
